The Instruction Decoder operates on 32\,bit instructions and generates all control signals needed for the operation of the CPU.

\paragraph{General Features: }
\begin{itemize}
  \item Decodes all instructions of the RV64I:
  \begin{itemize}
    \item opcode=3: I-type, loads
    \item opcode=19: I-type, immediate arithmetics, logic and shift
    \item opcode=23: U-type, auipc
    \item opcode=27: I-type, immediate arithmetics, logic and shift (word)
    \item opcode=35: S-type, stores
    \item opcode=51: R-type, RV32I
    \item opcode=55: U-type, lui
    \item opcode=59: R-type, RV64I
    \item opcode=99: B-type, all conditional branches
    \item opcode=103: I-type, jalr
    \item opcode=111: J-type, jal
  \end{itemize}
  \item generates the control signal for the ALU,
  \item ... for the D-Mem MUX,
  \item ... for the D-Mem read- and write enable,
  \item ... for the PC MUX,
  \item ... for the register A and B output MUXes,
  \item ... for the register file write enable,
  \item ... for the branch control,
  \item ... for the immediate generation selection,
  \item ... for the jump and branch control in the pipelined version. 
\end{itemize}

\paragraph{Synchronous Mode: }
\begin{itemize}
  \item -
\end{itemize}

\paragraph{Asynchronous Mode: }
\begin{itemize}
  \item -
\end{itemize}
